\section{Introduction}\label{sec:introduction}
A self-avoiding walk is a nearest-neighbor lattice path that visits no vertex twice. Its length is fixed combinatorially, but its spatial extent is random. The exponent $3/4$ expresses a predicted relation between those two quantities in two dimensions; its reciprocal $4/3$ expresses the corresponding mass dimension. A precise statement requires both a probability law and a choice of spatial statistic. This paper concerns the regular honeycomb lattice and gives such a statement for diameter, local mass, and covering numbers.

\subsection{The model and the meaning of the exponents}
Let $\mathbb H$ be the planar dual of the equilateral triangular tiling with triangle side length one. Fix a vertex $o$. Let $\mathcal W_n$ be the set of paths
\[
\gamma=(\gamma_0,\ldots,\gamma_n),\qquad \gamma_0=o,
\]
whose consecutive vertices are adjacent in $\mathbb H$ and whose vertices are all distinct. Write $c_n=|\mathcal W_n|$, and let $\Prob_n$ be the uniform law on $\mathcal W_n$. The endpoint $\gamma_n$ is free. We use the Euclidean embedding and set
\[
V(\gamma)=\{\gamma_0,\ldots,\gamma_n\},\quad
D(\gamma)=\diam V(\gamma),\quad
M_\gamma(z,s)=|V(\gamma)\cap\overline B(z,s)|.
\]
Let $\mathcal N_\gamma(s)$ be the least number of closed Euclidean balls of radius $s$ that cover $V(\gamma)$; their centers may lie anywhere in the plane.

\begin{theorem}[Fixed-length mass and covering exponents]\label{thm:main}\label{09_fast_paths:fixed-length}
For every $\delta>0$ and $k>0$, there are constants $C_{\delta,k}<\infty$ and $n_0(\delta,k)$ such that, for every integer $n\ge n_0(\delta,k)$, an event of $\Prob_n$-probability at least $1-C_{\delta,k}n^{-k}$ has all the following properties:
\begin{align}
 n^{3/4-\delta}&\le D(\gamma)\le n^{3/4+\delta},\label{eq:main-diameter}\\
 n^{-\delta}\min\{n,s^{4/3}\}
 &\le M_\gamma(z,s)\le n^{\delta}\min\{n,s^{4/3}\},\label{eq:main-mass}\\
 n^{-\delta}(1+ns^{-4/3})
 &\le \mathcal N_\gamma(s)\le n^{\delta}(1+ns^{-4/3}).\label{eq:main-cover}
\end{align}
The last two bounds hold simultaneously for every $z\in V(\gamma)$ and every real $s\in[1,n]$ (with $z$ relevant only to \eqref{eq:main-mass}).
\end{theorem}
The power slack absorbs constants and the difference between $n$ edges and $n+1$ vertices. In particular, the notation $n^{a+o(1)}$ below always has an explicitly stated probability interpretation; it does not assert a limiting multiplicative constant.

\begin{corollary}[Span in every lattice normal]\label{cor:directional-span}
For every \(\delta,k>0\), outside an event of \(\Prob_n\)-probability
\(O_{\delta,k}(n^{-k})\), the projection of \(V(\gamma)\) onto each
unit normal to a direction of triangular-lattice edges has span at least
\(n^{3/4-\delta}\). This holds for every sufficiently large integer \(n\).
\end{corollary}
This lower bound in each of the three normal directions gives more
information than a lower bound on the largest spatial extent.

\begin{corollary}[Covering dimension at polynomially separated scales]\label{cor:dimension}
For each fixed $\epsilon\in(0,3/4)$,
\[
\sup_{1\le s\le n^{3/4-\epsilon}}
\left|\frac{\log\mathcal N_\gamma(s)}{\log(D(\gamma)/s)}-\frac43\right|
\longrightarrow0
\quad\hbox{in }\Prob_n\hbox{-probability}.
\]
The statistic can be defined arbitrarily when $D(\gamma)\le n^{3/4-\epsilon}$.
\end{corollary}

\begin{corollary}[Radius of gyration and moments]\label{cor:gyration}
Let
\[
\bar\gamma=\frac1{n+1}\sum_{j=0}^n\gamma_j,
\qquad R_g(\gamma)^2=\frac1{n+1}\sum_{j=0}^n|\gamma_j-\bar\gamma|^2.
\]
Then $R_g=n^{3/4+o(1)}$ with the same arbitrary polynomial probability convention as in Theorem~\ref{thm:main}. For every fixed $p>0$,
\[
\E_n[D^p]=n^{3p/4+o(1)},\qquad
\E_n[R_g^p]=n^{3p/4+o(1)}.
\]
\end{corollary}
We prove the deductions from the geometric estimates, including the Corollaries, in Section~\ref{sec:fixed-length}. These statements concern the mass and covering behavior of the lattice trace as its length grows. The assertions do not include convergence to a continuum curve, a continuum Hausdorff dimension, a universality theorem on changing lattice, or a fixed-length lower bound on $|\gamma_n-o|$. A lower bound for the diameter does not imply that the two endpoints are far apart.


\subsection{Critical weights and the geometric argument}
The proof first works with the critical activity
\(\rho=(2+\sqrt2)^{-1/2}\). A path then carries weight \(\rho^\ell\),
where \(\ell\) counts edges for a vertex-rooted walk and visited vertices
for a path whose endpoints are ports, the midpoints of honeycomb edges. A \emph{mass} is a sum of
these weights before normalization. The distinction lets pieces in
adjacent open strips concatenate with exactly the product weight.
A bridge crosses a strip without visiting either boundary elsewhere.
It is irreducible if no intermediate lattice line parallel to the strip
boundaries is crossed exactly once; such a line would split it into two
bridges.

Exact lattice observables provide the mass crossing a strip, nesting
around a face, the length-square mass of polygons, and a uniform bound
for arcs joining two marked ports on a cylinder. We state their precise
positive outputs in Section~\ref{sec:inputs}. The analytic proofs belong
to the complete companion articles identified there. Here we prove the
geometric passage from those inputs to an unrestricted fixed-length trace.

There are three obstacles. A bridge used to modify a path must fit in a
corridor that avoids the remaining pieces; a strip mass or mean length
alone gives no such guarantee. We sew smaller bridges and recover the
sewing lines from the resulting path, controlling the multiplicity of
that operation. This produces localized first and second moments and,
by Cauchy--Schwarz, a subpower chance of seeing a long bridge. Independent
renewal trials amplify that chance. An exact conditioning identity then
transfers the estimate to a bridge of a specified height.

Next, one bridge estimate must control all subpaths of an arbitrary
walk. Ordered tests of adjacent renewal paths suppress both unusually
fast travel and excessive length in a narrow region. When the two tests
meet the same irreducible piece, its entire weight enters once. This
joint estimate is essential when the number of turns grows with the
scale. Finally, a separate estimate for many disjoint crossings controls
repeated visits to one ball. A subpath bound on its own would leave these
separate visits uncontrolled.

The estimates are proved for rooted walks of any length confined to
diameter at most a fixed multiple of \(H\). For every \(\tau>0\),
they simultaneously bound the length of every subwalk \(\sigma\) by
\(H^\tau(1+\operatorname{diam}\sigma)^{4/3}\), its diameter by
\(H^\tau(1+|\sigma|)^{3/4}\), and the number of visited vertices
in a ball of radius \(s\) about a visited vertex by \(H^\tau s^{4/3}\)
for \(1\le s\le H\). Here \(|\sigma|\) counts edges. For every
\(A>0\), the total critical mass of walks violating any of these
bounds is at most \(C_{\tau,A}H^{-A}\). This absolute estimate is
strong enough to survive summation over all the possible subpaths.

To pass to every integer length, submultiplicativity and the polynomial
strip lower bound give \(c_n\rho^n\ge1\). Taking \(H\) proportional
to \(n\) and dividing an exceptional mass by this fixed-length
partition function preserves arbitrary polynomial decay. The diameter,
local mass and covering conclusions then follow on one simultaneous event.
The elementary renewal process needed here is developed locally; no
change-of-ensemble theorem is assumed.

The proof follows these steps in order. Section~\ref{sec:full-box}
first bounds the mass of the outside pieces that a cut leaves behind.
Section~\ref{sec:08-sewing-renewal} proves the confined bridge length
moments and then constructs turning corridors for the later probes.
Section~\ref{sec:renewal} amplifies the moment bound at a specified
height. Section~\ref{sec:10-thermal-probes} then controls
ordered adjacent paths, including a shared long irreducible.
Sections~\ref{sec:fast} and~\ref{sec:11-slow-covering} use those estimates
to bound fast travel, excessive local length, and repeated visits.
Section~\ref{sec:fixed-length} performs the exact-length transfer and
proves the stated geometric consequences. The independent asynchronous
avoidance argument is retained in Appendix~\ref{sec:asynchronous}.


\subsection{Historical context}
The predicted exponent $3/4$ belongs to the two-dimensional polymer
theory developed through the dilute $O(n)$ model. Nienhuis's
solid-on-solid and Coulomb-gas analysis gives this exponent at $n=0$
\cite{Nienhuis1982}. These predictions concern more than the exponential
growth rate of the number of walks. Duminil-Copin and Smirnov proved the
exact honeycomb connective constant by introducing a parafermionic
observable and summing its local cancellation over a strip
\cite{DCS2012}. Their argument gives $c/T\le B_T\le1$ for the critical
strip-crossing mass; the sharper exponent $-1/4$ is recorded there as a
prediction.

Subsequent work sharpened the information about this mass.
Beaton, Bousquet-M\'elou, de Gier, Duminil-Copin and Guttmann proved
$B_T\to0$ in their study of surface adsorption \cite{BeatonEtAl2014}.
Glazman and Manolescu gave a shorter proof and a logarithmic subsequence
bound, and established invariance of the boundary two-point function
for a specified class of columnwise rhombic half-plane tilings
with integrable weights \cite{GlazmanManolescu2020}. More recently,
Krachun and Panagiotis proved
a polynomial upper bound on $B_T$ and a quantitative sub-ballistic
estimate: a uniform $n$-step honeycomb walk stays within distance
$C n/\log(n+1)$ of its origin except on an event of probability at most
$\exp(-c n^{2/3})$ \cite{KrachunPanagiotis2026}.

The analytic companions used here provide the further quantitative strip
and marked-polygon estimates needed for the geometric argument.

The bridge factorization follows Kesten's renewal method
\cite{Kesten1963}; see also Madras and Slade~\cite{MadrasSlade1993}
and the appendix of Lawler, Schramm and Werner~\cite{LSW2004}.
Dyhr, Gilbert, Kennedy, Lawler and
Passon~\cite[Proposition 2.10]{DyhrEtAl2011} identify the critical
free-endpoint strip law on the square lattice by conditioning the
infinite half-plane walk to have a cut at a prescribed height and
taking its prefix through the upper strip boundary.
We derive the renewal normalization and this finite-height identity
in the present port convention. For the final full-plane law at each
integer length, we use absolute exceptional masses and the
partition-function bound \(c_n\rho^n\ge1\).

Lawler, Schramm and Werner also proposed a continuum description through
\(\mathrm{SLE}_{8/3}\), whose Hausdorff dimension they identified as
\(4/3\) \cite[Proposition 2(ii)]{LSW2004}. Beffara subsequently proved
the general formula \(\min\{2,1+\kappa/8\}\) for the dimension of
\(\mathrm{SLE}_{\kappa}\) \cite{Beffara2008}.
These continuum results explain the predicted mass dimension;
a lattice scaling limit is neither an input nor an output of the
present proof.
